\begin{ex}
Khối lập phương có bao nhiêu mặt?
\shortans{6}
\loigiai{
Khối lập phương có 6 mặt.
}
\end{ex}

\begin{ex}
Khối lập phương có bao nhiêu đỉnh?
\shortans{8}
\loigiai{
Khối lập phương có 8 đỉnh.
}
\end{ex}

\begin{ex}
Khối lập phương có bao nhiêu cạnh?
\shortans{12}
\loigiai{
Khối lập phương có 12 cạnh.
}
\end{ex}

\begin{ex}
Các mặt của khối lập phương là hình gì? (Hình vuông ghi 1, hình chữ nhật ghi 2).
\shortans{1}
\loigiai{
Các mặt của khối lập phương đều là hình vuông.
}
\end{ex}

\begin{ex}
Khối hộp chữ nhật có bao nhiêu mặt?
\shortans{6}
\loigiai{
Khối hộp chữ nhật có 6 mặt.
}
\end{ex}

\begin{ex}
Khối hộp chữ nhật có bao nhiêu đỉnh?
\shortans{8}
\loigiai{
Khối hộp chữ nhật có 8 đỉnh.
}
\end{ex}

\begin{ex}
Khối hộp chữ nhật có bao nhiêu cạnh?
\shortans{12}
\loigiai{
Khối hộp chữ nhật có 12 cạnh.
}
\end{ex}

\begin{ex}
Tổng số mặt của 2 khối lập phương là bao nhiêu?
\shortans{12}
\loigiai{
$6 \times 2 = 12$ mặt.
}
\end{ex}

\begin{ex}
Tổng số cạnh của 2 khối lập phương là bao nhiêu?
\shortans{24}
\loigiai{
$12 \times 2 = 24$ cạnh.
}
\end{ex}

\begin{ex}
Tổng số đỉnh của 3 khối hộp chữ nhật là bao nhiêu?
\shortans{24}
\loigiai{
$8 \times 3 = 24$ đỉnh.
}
\end{ex}

\begin{ex}
Một con xúc xắc hình khối lập phương có các mặt được chấm từ 1 đến 6. Tổng số chấm trên 6 mặt là bao nhiêu?
\shortans{21}
\loigiai{
$1 + 2 + 3 + 4 + 5 + 6 = 21$.
}
\end{ex}

\begin{ex}
Một khối lập phương có cạnh dài 5 cm. Tổng độ dài tất cả 12 cạnh của khối lập phương đó là bao nhiêu xăng-ti-mét?
\shortans{60}
\loigiai{
$5 \times 12 = 60$ cm.
}
\end{ex}

\begin{ex}
Một khối lập phương có cạnh dài 3 cm. Diện tích một mặt của khối lập phương đó là bao nhiêu xăng-ti-mét vuông?
\shortans{9}
\loigiai{
$3 \times 3 = 9\text{ cm}^2$.
}
\end{ex}

\begin{ex}
Dùng các khối lập phương nhỏ cạnh 1 cm để xếp thành một khối lập phương lớn có cạnh dài 2 cm. Cần bao nhiêu khối lập phương nhỏ?
\shortans{8}
\loigiai{
$2 \times 2 \times 2 = 8$ khối.
}
\end{ex}

\begin{ex}
Dùng các khối lập phương nhỏ cạnh 1 cm xếp thành khối hộp chữ nhật dài 3 cm, rộng 2 cm, cao 1 cm. Cần bao nhiêu khối lập phương nhỏ?
\shortans{6}
\loigiai{
$3 \times 2 \times 1 = 6$ khối.
}
\end{ex}

\begin{ex}
Một hộp phấn hình khối lập phương có bao nhiêu cặp mặt đối diện song song và bằng nhau?
\shortans{3}
\loigiai{
Có 3 cặp mặt đối diện.
}
\end{ex}

\begin{ex}
Một thùng các-tông hình khối hộp chữ nhật có bao nhiêu đỉnh ở mặt đáy trên?
\shortans{4}
\loigiai{
Mỗi mặt đáy có 4 đỉnh.
}
\end{ex}

\begin{ex}
Một khối lập phương có tổng độ dài 12 cạnh là 48 cm. Độ dài mỗi cạnh của khối lập phương đó là bao nhiêu xăng-ti-mét?
\shortans{4}
\loigiai{
$48 : 12 = 4$ cm.
}
\end{ex}

\begin{ex}
Một khối lập phương có tổng độ dài 12 cạnh là 84 cm. Độ dài mỗi cạnh của khối lập phương đó là bao nhiêu xăng-ti-mét?
\shortans{7}
\loigiai{
$84 : 12 = 7$ cm.
}
\end{ex}

\begin{ex}
Xếp 4 khối lập phương nhỏ bằng nhau thành một hàng ngang. Khối hình tạo thành là khối gì? (Khối hộp chữ nhật ghi 1, khối lập phương ghi 2).
\shortans{1}
\loigiai{
Đó là khối hộp chữ nhật.
}
\end{ex}

\begin{ex}
Một bể cá hình hộp chữ nhật có 4 chân chống tại 4 đỉnh đáy. Bể cá đó có tất cả bao nhiêu đỉnh?
\shortans{8}
\loigiai{
Khối hộp chữ nhật có 8 đỉnh.
}
\end{ex}

\begin{ex}
Có bao nhiêu cạnh ở mặt đáy của khối hộp chữ nhật?
\shortans{4}
\loigiai{
Mặt đáy có 4 cạnh.
}
\end{ex}

\begin{ex}
Có bao nhiêu cạnh đứng nối hai mặt đáy của khối hộp chữ nhật?
\shortans{4}
\loigiai{
Có 4 cạnh bên đứng.
}
\end{ex}

\begin{ex}
Xếp các khối lập phương nhỏ cạnh 1 cm thành một khối lập phương có cạnh 3 cm. Cần bao nhiêu khối lập phương nhỏ?
\shortans{27}
\loigiai{
$3 \times 3 \times 3 = 27$ khối.
}
\end{ex}

\begin{ex}
Khối Rubik $3 \times 3 \times 3$ gồm bao nhiêu khối lập phương nhỏ?
\shortans{27}
\loigiai{
$3 \times 3 \times 3 = 27$ khối nhỏ.
}
\end{ex}
